\section{Additional Baseline Ablations}
\label{sec:baseline-ablations}

The main action-binding result compares UWM-JEPA-CF to a
parameter-matched LSTM-JEPA-CF configuration trained with the same
counterfactual target. We also ran LSTM-CF ablations that vary
recurrent state size and learning rate while leaving the
counterfactual-JEPA target and indicator evaluation unchanged.

\begin{table}[h]
\centering
\small
\caption{\textbf{LSTM-CF tuning confirmation at the original 5k-step
budget.} Values are mean $\pm$ std over three seeds. Target accuracy
reports the linear indicator readout on target-encoder latents; correct,
wrong, no-action, and shuffled report the same readout after blind
counterfactual rollout. Larger recurrent states improve target-latent
decodability in some configurations, but do not yield a stable
action-sensitive blind rollout.}
\label{tab:lstm-cf-confirm}
\begin{tabular}{lccccc}
\toprule
Configuration & Target & Correct & Wrong & No-action & Shuffled \\
\midrule
$h{=}12$, lr $3{\times}10^{-4}$ &
$0.645\pm0.220$ & $0.531\pm0.001$ & $0.531\pm0.001$ &
$0.530\pm0.000$ & $0.531\pm0.001$ \\
$h{=}32$, lr $1{\times}10^{-4}$ &
$0.622\pm0.087$ & $0.527\pm0.006$ & $0.525\pm0.009$ &
$0.525\pm0.009$ & $0.533\pm0.005$ \\
$h{=}63$, lr $1{\times}10^{-3}$ &
$0.867\pm0.071$ & $0.479\pm0.079$ & $0.511\pm0.041$ &
$0.504\pm0.045$ & $0.487\pm0.070$ \\
$h{=}255$, lr $1{\times}10^{-3}$ &
$0.568\pm0.010$ & $0.470\pm0.000$ & $0.470\pm0.000$ &
$0.470\pm0.000$ & $0.470\pm0.000$ \\
\bottomrule
\end{tabular}
\end{table}

The $h{=}63$ row is the clearest diagnostic. The target latents support
a high-accuracy linear indicator readout, but the blind rollout does
not preserve that readout under the correct action sequence. The
failure is therefore not simply absent target-state information in the
recurrent encoder. It appears when the counterfactual-JEPA predictor is
asked to carry that information forward without new observations.

\begin{table}[h]
\centering
\small
\caption{\textbf{Directly supervised recurrent positive control.}
Values are mean $\pm$ std over three seeds. The model receives the same
context observations and action sequence, but is trained directly for
the hidden-velocity indicator rather than through a JEPA latent target.}
\label{tab:lstm-supervised-control}
\begin{tabular}{lcc}
\toprule
Configuration & Correct actions & Wrong actions \\
\midrule
$h{=}12$  & $0.960\pm0.000$ & $0.867\pm0.017$ \\
$h{=}63$  & $0.955\pm0.015$ & $0.839\pm0.001$ \\
$h{=}255$ & $0.944\pm0.002$ & $0.848\pm0.011$ \\
\bottomrule
\end{tabular}
\end{table}

The supervised control shows that action-conditioned recurrent models
can solve the indicator when optimized directly for it. Together,
\cref{tab:lstm-cf-confirm,tab:lstm-supervised-control} locate the
main-comparison failure in the tested counterfactual-JEPA latent
rollout, not in the availability of action inputs or in recurrent
capacity alone.

We also evaluated a full-state linear probe that predicts both
continuous position and hidden velocity. Across the three UWM-JEPA-CF
seeds, blind-rollout full-state mean $R^2$ is $0.291\pm0.009$ and blind
velocity-sign accuracy is $0.734\pm0.007$. In the LSTM-CF confirmation
runs, target encoders can be decodable while blind rollout remains near
majority-class behavior; for example, the $h{=}63$ confirmation gives
target full-state mean $R^2=0.845\pm0.125$ but blind velocity-sign
accuracy $0.538\pm0.007$. This broader probe supports the same
interpretation as the binary indicator: the separation appears after
blind rollout, not simply at the target-encoder representation.
